\begin{figure}[htbp]
\centering
\begin{tikzpicture}[font=\small,line cap=round,line join=round]
  \begin{scope}[shift={(-3.35,2.15)}]
    \node[align=center,font=\bfseries\small] at (0,1.95) {Euclidean tangent geometry};
    \draw[mpblue!15,line width=.5pt] (-1.7,-1.1) grid[step=.55] (1.7,1.15);
    \filldraw[fill=mpblue!8,draw=mpblue,line width=1pt]
      (-1.1,-.65)--(1.1,-.65)--(.05,1.1)--cycle;
    \foreach \p in {(-1.1,-.65),(1.1,-.65),(.05,1.1)}\fill[mpblue] \p circle (2pt);
    \node[align=center,text width=4.5cm,text=mpgray] at (0,-1.65)
      {Shrink, then rescale\\about an interior point};
  \end{scope}
  \begin{scope}[shift={(3.35,2.15)}]
    \node[font=\bfseries\small] at (0,1.95) {Compact vertices};
    \draw[mpgray!65,line width=.6pt] (0,0) circle (1.35);
    \filldraw[fill=mpblue!8,draw=mpblue,line width=1pt]
      (90:.9)--(210:.9)--(330:.9)--cycle;
    \foreach \a in {90,210,330}\fill[mpblue] (\a:.9) circle (2pt);
    \node[text=mpgray] at (0,-1.65) {Vertices inside the Klein ball};
  \end{scope}
  \begin{scope}[shift={(-3.35,-2.35)}]
    \node[font=\bfseries\small] at (0,1.95) {Ideal vertices};
    \draw[mpgray!65,line width=.6pt] (0,0) circle (1.35);
    \filldraw[fill=mpteal!8,draw=mpteal,line width=1pt]
      (90:1.35)--(210:1.35)--(330:1.35)--cycle;
    \foreach \a in {90,210,330}\filldraw[fill=white,draw=mpteal,line width=1pt] (\a:1.35) circle (2.4pt);
    \node[text=mpgray] at (0,-1.65) {Vertices on the sphere at infinity};
  \end{scope}
  \begin{scope}[shift={(3.35,-2.35)}]
    \node[font=\bfseries\small] at (0,1.95) {Hyperideal vertices};
    \draw[mpgray!65,line width=.6pt] (0,0) circle (1.35);
    \draw[mpgray!60,dashed] (90:1.55)--(210:1.55)--(330:1.55)--cycle;
    \fill[mpteal!8]
      (-.21604075,1.17580645)--(-1.12629863,-.40080645)--(-.91025788,-.775)
      --(.91025788,-.775)--(1.12629863,-.40080645)--(.21604075,1.17580645)--cycle;
    \draw[mpblue,line width=1pt]
      (-.21604075,1.17580645)--(-1.12629863,-.40080645)
      (-.91025788,-.775)--(.91025788,-.775)
      (1.12629863,-.40080645)--(.21604075,1.17580645);
    \draw[mpochre,line width=1.4pt]
      (-.21604075,1.17580645)--(.21604075,1.17580645)
      (-1.12629863,-.40080645)--(-.91025788,-.775)
      (.91025788,-.775)--(1.12629863,-.40080645);
    \foreach \a in {90,210,330}\fill[mpgray] (\a:1.55) circle (1.6pt);
    \node[text=mpgray,align=center] at (0,-1.65)
      {Exterior vertices and\\polar truncation planes};
  \end{scope}
\end{tikzpicture}
\caption{Local pictures for the geometric regimes. The three Klein-model
drawings are planar sections illustrating the location of vertices;
the ochre segments are polar truncations in the hyperideal case.
The Euclidean picture represents a tangent-scale limit after shrinking
and rescaling. These drawings do not assert a single global deformation
through all four regimes with arbitrary prescribed data.}
\label{fig:mono-polyhedral-regimes}
\end{figure}
